\section{Agent 的最小定义：实体、环境与目标导向活动}

苏煜给出的 Agent 定义非常朴素：Agent 是一个有边界的实体，在某个外部环境中工作，并进行 goal-directed activities，也就是带有目标导向的活动。这个定义有意保持宽泛，因为动物、人、机器人、网页操作代理、coding agent 都可以落在这个框架里。关键不是它是否像人，而是它是否有可识别的边界、环境输入和目标驱动行为。

\begin{figure}[H]
\centering
\includegraphics[width=0.92\textwidth]{agent-loop.png}
\caption{Agent 的最小循环：环境、观察、推理、动作与记忆。}
\end{figure}

\begin{knowledgebox}{读图：Agent 循环应该怎么看}
图中最左侧是环境，可能是网页、桌面、代码库、数据库或机器人世界；上方 Observation 是感知到的状态；中间 Agent 根据目标、观察和记忆做推理；右侧 Action 改变环境；下方 Memory/World Model 把交互经验沉淀成可复用状态。它支持的结论是：Agent 不是单次回答，而是闭环系统。
\end{knowledgebox}

\begin{importantbox}{一个简化公式}
可以把 Agent 的一步动作写成：
$$
a_t = \pi(o_t, m_t, g)
$$
其中 $a_t$ 是第 $t$ 步动作；$o_t$ 是当前观察；$m_t$ 是记忆或世界模型；$g$ 是目标；$\pi$ 是策略或决策函数。这个公式不是访谈中的显式数学公式，而是对苏煜定义的教学化压缩。
\end{importantbox}

\subsection{为什么定义要宽}

如果定义太窄，Agent 很容易被误解成某种具体产品形态，例如网页代理、桌面代理或 coding agent。苏煜的定义把这些都看成同一类系统在不同环境中的实例。browser use、desktop use、mobile use、GUI、CLI、API、coding 都只是 means to an end；真正的目标是能在 digital world 中完成各种任务的 universal digital agent。

\begin{warningbox}{常见误解：Agent 不等于聊天框加插件}
聊天框调用一个工具，只是 Agent 的一个早期切片。真正的 Agent 要能感知状态、规划动作、执行动作、利用反馈，并在多步任务中保持目标。把 Agent 简化成“会调用 API 的 LLM”会低估 memory、world model、reliability 和 cost 的重要性。
\end{warningbox}

\subsection{本章小结}

Agent 的最小定义是：有边界的实体，在环境中基于目标采取行动。这个定义让我们能把早期符号系统、强化学习、semantic parsing、LLM 工具使用和现代 coding agent 放进同一条历史线里。

